\kbvcontinued ನ್ಯಾ, ವ್ಯಂಜನಾಕ್ಷರಗಳ ಹಿಂದಿನ ಅಕ್ಷರಗಳು, ಷಟ್ಪದಿಗಳಲ್ಲಿನ ಮೂರು ಮತ್ತು ಆರನೆಯ ಸಾಲುಗಳ ಅಂತ್ಯದಲ್ಲಿರುವ ಅಕ್ಷರಗಳು ‘ಗುರು’ಗಳೆನಿಸುವುವು. ಉದಾಹರಣೆಗೆ ಆ, ಈ, ಊ, ೠ, ಏ, ಐ, ಓ, ಔ; ಕಾ, ಕೀ, ಕೂ, ಕೄ, ಕೈ, ಕೇ, ಕೋ, ಕೌ, ಕಂ, ಕಃ, —ಗುರುಗಳು.\par

\indent ಪದ್ಯದ ಸಾಲಿನಲ್ಲಿ ಬರುವ ಅಕ್ಷರಗಳ ಮೇಲೆ ‘\guru{}’ ಚಿಹ್ನೆ ಹಾಕಿ ಗುರುಗಳನ್ನು ಗುರ್ತಿಸುವುದು ವಾಡಿಕೆ.\par

\vspace{3pt plus 2pt}\noindent\hspace*{1.2em}\kbvlabel{ಉದಾ}\par\vspace{1pt}

\vspace{2.5pt plus 3pt minus 0.5pt}\noindent\hspace*{1.2em}\kbvnum{(1)}\hspace{0.55em}ಓಷಧ\par

\vspace{2.5pt plus 3pt minus 0.5pt}\noindent\hspace*{1.2em}\kbvnum{(2)}\hspace{0.55em}ಕಾಗದ\par

\vspace{2.5pt plus 3pt minus 0.5pt}\noindent\hspace*{1.2em}\kbvnum{(3)}\hspace{0.55em}ಮಂಗಳ\par

\vspace{2.5pt plus 3pt minus 0.5pt}\noindent\hspace*{1.2em}\kbvnum{(4)}\hspace{0.55em}ದುಃಖ\par

\vspace{2pt plus 2pt}\noindent\hspace*{1.5em}(5) ಉತ್ಸವ\hfill (6) ಈಶ್ವರ\hfill (7) ಕೊರಳ್\hspace*{1.5em}\par\vspace{2pt plus 2pt}

\vspace{2.5pt plus 3pt minus 0.5pt}\noindent\hspace*{1.2em}\kbvnum{(1)}\hspace{0.55em}\guru{}\laghu{}\laghu{} (ಇಲ್ಲಿ ‘ಓ’ ಎಂಬುದು ದೀರ್ಘಸ್ವರ)\par

\vspace{2.5pt plus 3pt minus 0.5pt}\noindent\hspace*{1.2em}\kbvnum{(2)}\hspace{0.55em}\guru{}\laghu{}\laghu{} (ಇಲ್ಲಿ ‘ಕಾ’ ಎಂಬುದು ದೀರ್ಘಸ್ವರದಿಂದ ಕೂಡಿದ ಗುಣಿತಾಕ್ಷರ)\par

\vspace{2.5pt plus 3pt minus 0.5pt}\noindent\hspace*{1.2em}\kbvnum{(3)}\hspace{0.55em}\guru{}\laghu{}\laghu{} (ಇಲ್ಲಿ ‘ಮಂ’ ಎಂಬುದು ಅನುಸ್ವಾರ)\par

\vspace{2.5pt plus 3pt minus 0.5pt}\noindent\hspace*{1.2em}\kbvnum{(4)}\hspace{0.55em}\guru{} (ಇಲ್ಲಿ ‘ದುಃ’ ಎಂಬುದು ವಿಸರ್ಗ)\par

\vspace{2.5pt plus 3pt minus 0.5pt}\noindent\hspace*{1.2em}\kbvnum{(5)}\hspace{0.55em}\guru{}\laghu{}\laghu{} (ಇಲ್ಲಿ ‘ಉ’ ಎಂಬುದು ಒತ್ತಕ್ಷರದ ಹಿಂದಿನ ಅಕ್ಷರ)\par

\vspace{2.5pt plus 3pt minus 0.5pt}\noindent\hspace*{1.2em}\kbvnum{(6)}\hspace{0.55em}\guru{}\laghu{}\laghu{} (ಇಲ್ಲಿ ‘ಈ’ ಎಂಬುದು ಒತ್ತಕ್ಷರದ ಹಿಂದಿನ ಅಕ್ಷರ ಹಾಗೂ ದೀರ್ಘಸ್ವರ ಆಗಿರುವುದು)\par

\vspace{2.5pt plus 3pt minus 0.5pt}\noindent\hspace*{1.2em}\kbvnum{(7)}\hspace{0.55em}\laghu{}\guru{} (ಇಲ್ಲಿ ‘ರ’ ಎಂಬುದು ‘ಳ್’ ವ್ಯಂಜನಾಕ್ಷರದ ಹಿಂದಿನ ಅಕ್ಷರ)\par

\indent ಸೂಚನೆ: ವ್ಯಂಜನಾಕ್ಷರಗಳನ್ನು (ಉದಾ: ಕ್, ಟ್, ನ್, ಭ್, ಣ್ ಇತ್ಯಾದಿ) ಲಘುಗಳೆಂದಾಗಲೀ, ಗುರುಗಳೆಂದಾಗಲೀ ಪರಿಗಣಿಸಲಾಗದು ಆದುದರಿಂದ ಅವುಗಳಿಗೆ ಯಾವ ಚಿಹ್ನೆಯನ್ನೂ ಹಾಕಬಾರದು.\par
